\chapter{Công trình liên quan}
\label{chap:lienquan}

Chương này khảo sát các nghiên cứu đã công bố (mục~\ref{sec:nghiencuulienquan}) và các nền tảng đã triển khai trong nước, quốc tế (mục~\ref{sec:hethonglienquan}); mục~\ref{sec:visaochualam} phân tích vì sao chưa nền tảng nào đi theo hướng của Farmery, và mục~\ref{subsec:researchgap} định vị khoảng trống đồ án hướng tới lấp đầy.

\section{Các nghiên cứu liên quan}
\label{sec:nghiencuulienquan}

Bốn công trình dưới đây mỗi công trình soi vào một khía cạnh của bài toán:
\begin{itemize}
    \item \textbf{Nền tảng thương mại điện tử và nông hộ nhỏ.} Morepje và cộng sự (2024) tổng quan giai đoạn 2010--2024 tại châu Phi hạ Sahara: nền tảng số cải thiện tiếp cận thị trường, giúp bán giá tốt hơn và giảm tổn thất sau thu hoạch, nhưng còn hai rào cản dai dẳng là hạ tầng số thiếu và trình độ số của nông dân thấp \parencite{morepje2024ecommerce}. \emph{Liên hệ:} hai rào cản này trùng khó khăn của nhóm nhà cung cấp (mục~\ref{sec:doituongnguoidung}) và là căn cứ cho yêu cầu giao diện tối giản (NFR6); các nền tảng trong nghiên cứu thuần trung gian nên nông dân vẫn tự đóng gói và giao đơn nhỏ --- đúng chỗ luồng bán lẻ 1P nhắm tới.
    \item \textbf{Chuẩn định danh cho truy xuất nguồn gốc.} Keogh và cộng sự (2020) kết luận chuẩn GS1 kết hợp blockchain, dù cần thiết, \emph{không tự giải quyết} bất cân xứng thông tin vì chuẩn không bảo đảm dữ liệu đưa vào là đúng \parencite{keogh2020blockchain}. \emph{Liên hệ:} lập luận nền cho việc dùng nguyên lý GS1 trên cơ sở dữ liệu tập trung (mục~\ref{subsec:gs1blockchain}) và đặt kiểm duyệt chứng nhận của con người làm lớp bảo đảm dữ liệu đầu vào.
    \item \textbf{Học máy trong dự báo giá nông sản.} Tran và cộng sự (2023) ghi nhận tiềm năng nhưng chỉ ra hạn chế về chất lượng và độ dài chuỗi dữ liệu lịch sử \parencite{tran2023pricepredict}. \emph{Liên hệ:} căn cứ để ưu tiên mô hình chuỗi thời gian có thành phần mùa vụ thay vì học sâu, đánh giá offline trên dữ liệu giá công khai (mục~\ref{sec:giaiphapcongnghe}).
    \item \textbf{Đề xuất tích hợp gần nhất.} Farm2Market (2025) kết hợp blockchain, ký quỹ smart contract, AI dự báo giá (độ chính xác khoảng 82\%) và QR truy xuất \parencite{farm2market_2025}. \emph{Liên hệ:} mới dừng ở đề xuất, chưa có dữ liệu vận hành và không phục vụ song song bán lẻ và bán sỉ; con số 82\% là mốc để Farmery đặt ngưỡng dự báo giá thấp hơn một chút (Phụ lục~E).
\end{itemize}

\section{Các hệ thống và nền tảng liên quan}
\label{sec:hethonglienquan}

\subsection{Các sàn thương mại điện tử trong nước}
\label{subsec:sosanhtrongnuoc}

Nhóm sàn tổng hợp lớn nhất là Shopee (56\% giá trị giao dịch của nhóm bốn sàn lớn trong chín tháng đầu 2025) và Lazada (thu hẹp còn 3\%, từ 6\% năm 2024) \parencite{vneconomy_2025_ecommerce}; cả hai có ngành thực phẩm nhưng không chuyên biệt hóa cho nông sản. Nhóm chuyên biệt gồm FoodMap, sàn tư nhân bán lẻ và bán sỉ với hơn 300 nông trại và QR truy xuất \parencite{foodmap}, và các sàn của bưu chính nhà nước: Postmart.vn và Vỏ Sò từng giúp tiêu thụ gần 1.200 tấn nông sản đợt dịch 2021 \parencite{postmart_voso_2021}, nhưng Vỏ Sò ngừng hoạt động đầu 2023, Postmart.vn đổi thành Buudien.vn \parencite{buudien_2024}, và đến 12/2024 cả hai ra mắt sàn mới theo hướng bán sỉ, xuất khẩu (nongsan.buudien.vn, VIPO Mall) \parencite{vipo_nongsan_2024}. Bảng~\ref{tab:sosanhtrongnuoc} đối chiếu các nền tảng này với Farmery.

\begingroup
\tblsetup\tblzebra
\begin{longtable}{L{2.53cm} L{3.13cm} L{4.44cm} L{4.44cm}}
\caption{So sánh các sàn thương mại điện tử trong nước liên quan đến nông sản}
\label{tab:sosanhtrongnuoc} \\
\toprule
\rowcolor{tblheadbg}
\textbf{Nền tảng} & \textbf{Mô hình} & \textbf{Điểm mạnh} & \textbf{Hạn chế với nông sản} \\
\midrule
\endfirsthead
\toprule
\rowcolor{tblheadbg}
\textbf{Nền tảng} & \textbf{Mô hình} & \textbf{Điểm mạnh} & \textbf{Hạn chế với nông sản} \\
\midrule
\endhead
\bottomrule
\endfoot
Shopee, Lazada & Sàn tổng hợp, B2C, 3P & Thị phần lớn; logistics và thanh toán trưởng thành & Không chuyên nông sản; không xác minh chứng nhận, không truy xuất theo lô, không giá sỉ theo MOQ \\
FoodMap & Chuyên nông sản, B2C và B2B, chủ yếu 3P & QR truy xuất; hơn 300 nông trại; đã vận hành thực tế & Quy mô nhỏ; chưa rõ giá bậc theo MOQ; chưa có AI hỗ trợ ra quyết định \\
Buudien.vn, VIPO Mall & Chuyên nông sản, B2B và xuất khẩu, 3P & Hạ tầng bưu chính, Viettel Post; tập trung xuất khẩu & Thu hẹp phạm vi; không có kênh bán lẻ trên cùng nền tảng \\
\textbf{Farmery} & Bán sỉ 3P kết hợp bán lẻ 1P & Hai kênh trên cùng nguồn cung; truy xuất gắn nhật ký canh tác; kiểm duyệt chứng nhận ba tầng; ba mô-đun AI & Chưa kiểm chứng ở quy mô thật; cần vốn lưu động, chịu rủi ro tồn kho ở bán lẻ \\
\end{longtable}
\endgroup

\subsection{Các nền tảng quốc tế}
\label{subsec:sosanhquocte}

Ba nền tảng quốc tế minh họa các hướng số hóa nông nghiệp đã kiểm chứng thương mại, mỗi nền tảng đều có giới hạn mà Farmery nhắm khắc phục. \textbf{e-Choupal} (ITC, Ấn Độ, từ 2000) dùng kiosk internet tại làng để bỏ tầng trung gian, tăng thu nhập nông dân tới 50\% ở quy mô khoảng 4 triệu nông dân \parencite{business_standard_2020_echoupal}, nhưng chỉ phục vụ chuỗi nội bộ của ITC, không phải marketplace mở. \textbf{Pinduoduo} (Trung Quốc) dùng mua chung nối trực tiếp hơn 12 triệu nông dân với người tiêu dùng \parencite{nasdaq_2021_pinduoduo}, nhưng thuần B2C, không có bán sỉ theo MOQ và không nhấn mạnh truy xuất. \textbf{Ninjacart} (Ấn Độ) dùng học máy dự báo nhu cầu, phục vụ hơn 17.000 điểm bán lẻ trong 12 giờ và trả tiền nông dân trong 24 giờ \parencite{ninjacart}, nhưng thuần B2B và chi phí kho lạnh cao.

\section{Vì sao các nền tảng hiện có chưa đi theo hướng này}
\label{sec:visaochualam}

Nếu mô hình kết hợp bán sỉ qua sàn với bán lẻ tự thu mua hợp lý, vì sao các nền tảng đủ nguồn lực lại chưa làm? Có bốn nguyên nhân:

\begin{enumerate}
    \item[\textbf{N1.}] \textbf{Mô hình phí của sàn tổng hợp không khớp hàng tươi:} hoa hồng và quảng cáo giả định chi phí xử lý mỗi đơn gần bằng không, trong khi nông sản tươi giá trị đơn thấp, tỷ lệ khiếu nại và hoàn hàng cao, mỗi khiếu nại cần người xử lý.
    \item[\textbf{N2.}] \textbf{Tự sở hữu hàng đòi hỏi vốn lưu động và hạ tầng bảo quản:} nền tảng bỏ tiền trước và chịu rủi ro hàng hỏng, nhất là khi kho lạnh rau quả tươi còn mỏng (Chương~\ref{chap:gioithieu}).
    \item[\textbf{N3.}] \textbf{Né rủi ro tồn kho bằng thiết kế khác:} Pinduoduo gom đủ cầu bằng mua chung trước khi thu hoạch; e-Choupal chỉ số hóa thông tin và thu mua.
    \item[\textbf{N4.}] \textbf{Chuyên môn hóa một phía thị trường:} Ninjacart chỉ phục vụ doanh nghiệp; FoodMap tiếp cận người tiêu dùng nhưng chưa có đàm phán giá sỉ theo MOQ.
\end{enumerate}

Farmery chấp nhận ba đánh đổi mà các nền tảng trên đã tránh: cần vốn lưu động cho thu mua, chịu lỗ trực tiếp với lô không bán hết, và vận hành hai mô hình kinh doanh trên một hệ thống. Điều làm đánh đổi này chấp nhận được là phạm vi thí điểm hẹp về địa bàn và chủng loại hàng (mục~\ref{sec:phamvi}).

\section{Khoảng trống nghiên cứu}
\label{subsec:researchgap}

Bảng~\ref{tab:researchgap} đối chiếu các nền tảng theo cùng một bộ tiêu chí. Ký hiệu: \checkmark{} đáp ứng đầy đủ, $\circ$ một phần hoặc không rõ, $\times$ không đáp ứng.

\begingroup
\tblsetup\tblzebra
\begin{table}[H]
\centering
\small
\caption{Ma trận khoảng trống nghiên cứu: Farmery so với các nền tảng đối chứng}
\label{tab:researchgap}
\begin{tabular}{L{3.49cm} >{\centering\arraybackslash}p{1.46cm} >{\centering\arraybackslash}p{1.46cm} >{\centering\arraybackslash}p{1.65cm} >{\centering\arraybackslash}p{1.55cm} >{\centering\arraybackslash}p{1.84cm} >{\centering\arraybackslash}p{2.04cm}}
\toprule
\rowcolor{tblheadbg}
\textbf{Tiêu chí} & \textbf{Shopee} & \textbf{Lazada} & \textbf{Food\-Map} & \textbf{VIPO Mall} & \textbf{Farm2\-Market} & \textbf{Farmery} \\
\midrule
Song song bán lẻ và bán sỉ theo MOQ & $\times$ & $\times$ & $\circ$ một phần & $\times$ chỉ B2B & $\times$ không nêu & \checkmark \\
Mô hình sở hữu hàng hóa & $\times$ chỉ 3P & $\times$ chỉ 3P & $\circ$ chủ yếu 3P & $\times$ chỉ 3P & $\times$ không nêu & \checkmark{} 3P và 1P \\
Truy xuất gắn nhật ký canh tác & $\times$ & $\times$ & $\circ$ QR cơ bản & $\circ$ không rõ & \checkmark{} blockchain & \checkmark{} QR và nhật ký \\
Xác thực chứng nhận GAP & $\times$ & $\times$ & $\circ$ không rõ & $\circ$ không rõ & $\times$ không nêu & \checkmark{} hai cấp \\
AI hỗ trợ ra quyết định & $\circ$ không cho nông sản & $\circ$ không cho nông sản & $\times$ & $\times$ & \checkmark{} dự báo giá & \checkmark{} ba mô-đun \\
Đã kiểm chứng thương mại & \checkmark{} 56\% GMV & $\circ$ thu hẹp, 3\% & \checkmark & \checkmark & $\times$ mới nghiên cứu & $\circ$ đang xây dựng \\
\bottomrule
\end{tabular}
\end{table}
\endgroup

Khoảng trống Farmery hướng tới gồm ba phần: truy xuất gắn nhật ký canh tác theo nguyên lý GS1 (mục~\ref{subsec:gs1blockchain}); phục vụ song song hai luồng trên cùng nguồn cung --- bán sỉ theo MOQ do nhà cung cấp tự bán và bán lẻ do nền tảng thu mua rồi bán lại; và tích hợp các mô-đun AI hỗ trợ ra quyết định. Tổ hợp này chưa được nền tảng đối chứng hiện thực đủ cùng lúc, vì những lý do ở mục~\ref{sec:visaochualam}.
